\section{Fribourg, Écône, and a Lawful Beginning}

The Society did not begin as a clandestine episcopal project. Seminarians dissatisfied with changing formation approached Lefebvre, and a house of formation opened at Fribourg with ecclesiastical permission. A letter of 18 August 1970 from Bishop François Charrière recorded his permission of 6 June 1969 for the Fribourg house. On 1 November 1970 he erected the Priestly Society of Saint Pius X in the Diocese of Lausanne, Geneva and Fribourg as a \emph{pia unio}. The same decree approved its statutes for six years \emph{ad experimentum}, with tacit renewal and the possibility of a later definitive erection. The decree named Fribourg as its seat and placed the work within the diocesan framework then available under the 1917 Code.

Écône, a former religious house in Valais, became the seminary. Its location in the Diocese of Sion and the Society's erection in another Swiss diocese later complicated public shorthand about ``approval.'' The original act was real but limited: diocesan, with experimentally approved statutes, and dependent on ecclesiastical law. It was not a papal erection, a permanent privilege, or advance authorization for a worldwide ministry independent of local ordinaries.

On 18 February 1971 Cardinal John Wright, prefect of the Congregation for the Clergy, wrote a letter praising the statutes' capacity to foster priestly life. The letter is important evidence that the young work was known favorably in Rome. It is not itself a decree erecting the Society as an institute of pontifical right. SSPX retrospectives sometimes place the letter under a heading of Roman approval; the juridically safer description preserves the distinction between favorable correspondence and a constitutive act.

\statusframe{Charrière's decree of 1 November 1970 and Wright's letter of 18 February 1971.}{The first erected a diocesan pious union and separately approved its statutes for six years experimentally, with tacit renewal; the second gave favorable Roman testimony to those statutes.}{Neither instrument supplied a permanent universal mission, and the cardinal's letter did not convert the diocesan body into one of pontifical right.}

\subsection{A seminary becomes a movement}

Écône offered a complete institutional package: common life, spiritual discipline, Thomistic studies, the older Roman rite, and a founder able to ordain. Vocations arrived from several countries. The first priestly ordinations and new houses turned a formation initiative into an apostolic network. The emerging structure joined common formation and priestly life to works related to them. It later supported houses, schools, chapels, retreats, publications, and affiliated communities.

Growth also made the unresolved questions unavoidable. A bishop can form seminarians only within a system that determines incardination, dimissorial letters, faculties, assignments, and relations with local ordinaries. The more priests Écône produced for work across dioceses, the less the dispute could remain a pedagogical disagreement inside one approved experiment. The central historical transition was therefore not from ``legal'' to ``illegal'' on a single self-explanatory day. It was from a bounded diocesan work to a transnational program whose founder and ecclesiastical superiors no longer agreed on the conditions of its lawful continuation.

\subsection{The name and patron}

The Society took Saint Pius X as patron, signaling a program of priestly sanctification, doctrinal resistance to modernism, and attachment to inherited worship. The name also framed later memory: supporters could present continuity with a canonized reforming pope, while opponents saw selective appeal to an earlier papacy against subsequent councils and popes. A historical account need not decide the theological validity of either symbolism to recognize its organizing force. The name made the Society's claim larger than preservation of one seminary: it presented its work as a defense of Catholic Tradition for the whole Church.
